\section{Introduzione}

Questo report descrive il lavoro svolto nell'ambito del progetto del corso ISW2,
articolato in quattro milestone progressive. L'obiettivo complessivo è investigare la
relazione tra code smell, bugginess e possibilità di refactoring automatico in un progetto
software open source reale.

Il progetto analizzato è \textbf{Apache OpenJPA}, un framework open source per la
persistenza dei dati in Java, sviluppato nell'ambito della Apache Software Foundation.
I dati sui ticket di bug sono stati recuperati tramite le API REST di Jira
(\texttt{issues.apache.org}), mentre il codice sorgente è stato estratto dal repository
Git ufficiale del progetto.

Le quattro milestone affrontate sono le seguenti:

\begin{enumerate}
    \item \textbf{Milestone 1 -- Costruzione del dataset}: raccolta e sanitizzazione dei
    ticket Jira, estrazione degli snapshot di classi Java per release, calcolo delle
    metriche, etichettatura della bugginess.

    \item \textbf{Milestone 2 -- Identificazione di modelli accurati}: feature selection, data balancing, addestramento e
    valutazione di classificatori per la predizione della bugginess a partire dalle
    metriche calcolate nella Milestone 1.

    \item \textbf{Milestone 3 -- Analisi dell'impatto dei code smell}: investigazione di
    quante classi buggy sarebbero potute essere prevenute in assenza di code smell.

    \item \textbf{Milestone 4 -- Refactoring automatico}: valutazione della possibilità
    di eliminare automaticamente i code smell tramite refactoring guidato, sotto diverse
    condizioni di disponibilità di test.
\end{enumerate}
